\documentclass[twocolumn]{aastex631}
\begin{document}
\title{A residual reionization ionizing-photon-budget shortfall for star-forming galaxies at z\~{}11 (jwsthigh systematic corner)}
\author{NebulaMind Lab (autonomous pipeline)}
\affiliation{NebulaMind Lab, \url{https://lab.nebulamind.net}}
\begin{abstract}
Reionization ionizing-photon-budget reconciliation at z\~{}11: star-forming galaxies require f\_esc=0.466 (+0.440/-0.224) to close the budget (Madau-Dickinson SFRD, log xi\_ion=25.5+/-0.15, clumping C=2-5, JWST-SFRD tail), versus indirect-proxy-inferred f\_esc=0.062 (+0.110/-0.039) from LzLCS O32/beta calibrations. Median delta(required-inferred)=+0.376 dex-frac (16-84\%: +0.138 to +0.818); 95\% of the systematic MC shows a shortfall. Conclusion: a genuine SHORTFALL remains, and the sign holds under both O32 and beta calibrations. The result is bounded by the xi\_ion x clumping x proxy-calibration systematic, not by statistics. Generated autonomously from public data (jwst) via the NebulaMind Lab runner: a bounded, reproducible, descriptive study.
\end{abstract}
\section{Introduction}
Introduction: Recent studies have highlighted a potential crisis in our understanding of the reionization process, with concerns that the observed star-forming galaxies may not produce enough ionizing photons to account for the rapid reionization suggested by observations [Muñoz2024]. This has led to increased scrutiny of the assumptions and models used to estimate the ionizing photon budget. As noted by Davies et al. (2021), the demands on ionizing sources are particularly challenging during the absorption-dominated phase of reionization. In response, researchers have sought to refine their understanding of the ionizing photon production efficiency through improved calibrations [Park2022] and assessments of the galaxy ionizing photon budget at various redshifts [Duncan2015].
\section{Data and method}
Data and method: To address this issue, we adopted a literature-anchored approach, relying on established values from previous studies. Specifically, we used the cosmic star formation rate density (SFRD) derived by Madau \& Dickinson (2014), along with published calibrations for the ionizing photon production efficiency (xi\_ion) and escape fraction (f\_esc) proxies [Chisholm+22, Flury+22; Simmonds+24]. Our method focused on reconciling these values to determine whether star-forming galaxies can account for the required ionizing photons during reionization at z\~{}11.
\section{Result}
Result: Based on our calculations, we found that star-forming galaxies would need an escape fraction of f\_esc = 0.466 (+0.440/-0.224) to close the ionizing photon budget, assuming a Madau-Dickinson SFRD and log xi\_ion = 25.5 ± 0.15, with clumping factor C ranging from 2 to 5. However, indirect-proxy-inferred values of f\_esc derived from LzLCS O32/beta calibrations yield a significantly lower estimate of f\_esc = 0.062 (+0.110/-0.039). This discrepancy results in a median shortfall of +0.376 dex-frac (ranging from +0.138 to +0.818), with 95\% of our systematic Monte Carlo simulations indicating a genuine shortfall.
\begin{figure}\centering\includegraphics[width=\columnwidth]{result.png}\caption{A residual reionization ionizing-photon-budget shortfall for star-forming galaxies at z\~{}11 (jwsthigh systematic corner)}\end{figure}
\section{Caveats}
Caveats: It is essential to acknowledge the limitations of our approach, which relies on automated single-selection and uncalibrated measurements. The accuracy of our results depends heavily on the assumptions underlying the adopted calibrations and models, including the Madau-Dickinson SFRD and O32/beta proxy relationships. Additionally, systematic uncertainties in the clumping factor and ionizing photon production efficiency may further impact our conclusions. Therefore, while our study highlights a potential shortfall in the ionizing photon budget, it is crucial to recognize that these findings are contingent upon the validity of the input parameters and calibrations used. Further research is needed to refine these estimates and address the inherent uncertainties in this line of inquiry.
\section*{References}
\footnotesize
[Muoz2024] Muñoz et al. (2024). Reionization after JWST: a photon budget crisis?. bibcode 2024MNRAS.535L..37M \par [Davies2021] Davies et al. (2021). The Predicament of Absorption-dominated Reionization: Increased Demands on Ionizing Sources. bibcode 2021ApJ...918L..35D \par [Park2022] Park et al. (2022). Calibrating excursion set reionization models to approximately conserve ionizing photons. bibcode 2022MNRAS.517..192P \par [Duncan2015] Duncan et al. (2015). Powering reionization: assessing the galaxy ionizing photon budget at z < 10. bibcode 2015MNRAS.451.2030D \par [Madau2017] Madau (2017). Cosmic Reionization after Planck and before JWST: An Analytic Approach. bibcode 2017ApJ...851...50M

\end{document}
